% ============================================================
% APPENDICES
% ============================================================
\appendix
\section{MedScript Grammar}

The grammar below is the one implemented by the recursive-descent parser. Terminals are in quotes; \code{ID},
\code{NUM}, \code{STR} and \code{UNIT} are token classes.

\begin{lstlisting}[language={},numbers=none,basicstyle=\ttfamily\footnotesize]
program    = { decl } ;
decl       = inputDecl | constDecl | funcDecl | stmt ;
inputDecl  = "input" ID ":" type [ "phi" ]
             [ "range" "(" signedNum "," signedNum ")" ] ";" ;
constDecl  = "const" ID ":" type "=" expr ";" ;
funcDecl   = "func" ID "(" [ param { "," param } ] ")" "->" type block ;
param      = ID ":" type ;
type       = ( "int" | "float" | "bool" | "string" | unitType ) [ "?" ] ;
unitType   = unitTerm { ( "*" | "/" ) unitTerm } ;
unitTerm   = UNIT [ "^" NUM ] ;

stmt       = letStmt | ifStmt | whileStmt | returnStmt | assignStmt
           | exprStmt | block ;
letStmt    = "let" ID [ ":" type ] [ "=" expr ] ";" ;
assignStmt = ID "=" expr ";" ;
ifStmt     = "if" "(" expr ")" block [ "else" ( ifStmt | block ) ] ;
whileStmt  = "while" "(" expr ")" block ;
returnStmt = "return" [ expr ] ";" ;
exprStmt   = expr ";" ;
block      = "{" { stmt } "}" ;

expr       = orExpr ;
orExpr     = andExpr { "||" andExpr } ;
andExpr    = eqExpr  { "&&" eqExpr } ;
eqExpr     = relExpr { ( "==" | "!=" ) relExpr } ;
relExpr    = addExpr { ( "<" | ">" | "<=" | ">=" ) addExpr } ;
addExpr    = mulExpr { ( "+" | "-" ) mulExpr } ;
mulExpr    = unary   { ( "*" | "/" | "%" ) unary } ;
unary      = ( "-" | "!" ) unary | primary ;
primary    = NUM [ unitType ] | STR | "true" | "false"
           | ID [ "(" [ expr { "," expr } ] ")" ] | "(" expr ")" ;
\end{lstlisting}

A unit suffix after a number extends over \code{/} only when the next token is a unit, so \code{15 mg/kg} is a single
literal while \code{15 mg / weight} is a division.

\section{Diagnostic Catalogue}

\begin{longtable}{llllp{6.3cm}}
\toprule
\textbf{Code} & \textbf{Phase} & \textbf{Severity} & \textbf{Family} & \textbf{Meaning}\\
\midrule
\endhead
LEX001 & lexer & error & lexical & illegal character\\
LEX002 & lexer & error & lexical & malformed number\\
LEX003 & lexer & error & lexical & unterminated string literal\\
SYN001 & parser & error & syntax & expected token not found\\
SYN002 & parser & error & syntax & unexpected token in expression\\
SYN003 & parser & error & syntax & unexpected end of file\\
SEM001 & semantic & error & name & undeclared identifier or function\\
SEM002 & semantic & error & name & redeclaration\\
SEM003 & semantic & error & name & use before initialisation\\
SEM009 & semantic & error & name & assignment to an input or constant\\
SEM004 & semantic & error & type & type mismatch\\
SEM006 & semantic & error & type & wrong argument count or type\\
SEM008 & semantic & error & type & condition is not boolean\\
SEM005 & semantic & error & unit & unit-of-measure mismatch\\
DOM001 & semantic & error & range & implausible value or threshold outside declared range\\
DOM002 & semantic & error & phi & PHI reaches an output sink\\
DOM003 & semantic & error & missing & optional value used without presence check\\
SEM007 & semantic & error & logic & missing return on some path\\
DOM004 & semantic & error & logic & division by constant zero\\
DOM005 & semantic & error & logic & divisor range includes zero\\
W001 & semantic & warning & logic & unreachable code\\
W002 & semantic & warning & valid & unused variable\\
W003 & semantic & warning & logic & constant condition\\
\bottomrule
\end{longtable}

\section{Fault-Injection Operators}

Table~\ref{tab:ops} lists the operators together with the number of times each was attempted while building the
corpus and how often the resulting fault was kept. \emph{Tried} counts every attempt, including those where the
operator was not applicable to the sampled program (for example, removing an \code{anonymize} call from a program
that has none); \emph{Kept} counts attempts that produced a program in which the fault manifested. The percentage is
therefore a mixture of applicability and manifestation, and it is not a measure of compiler accuracy.

{\small
\begin{longtable}{llp{5.2cm}rrr}
\caption{Fault-injection operators and yield during corpus construction.}\label{tab:ops}\\
\toprule
\textbf{Family} & \textbf{Operator} & \textbf{Fault injected} & \textbf{Tried} & \textbf{Kept} & \textbf{\%}\\
\midrule
\endhead
lexical & \texttt{lex\_illegal} & insert an illegal character & 332 & 332 & 100\\
 & \texttt{lex\_number} & malform a numeric literal & 311 & 307 & 99\\
 & \texttt{lex\_string} & drop a closing quote & 362 & 361 & 100\\
\midrule
syntax & \texttt{syn\_semicolon} & delete a semicolon & 170 & 170 & 100\\
 & \texttt{syn\_paren} & delete a closing parenthesis & 173 & 173 & 100\\
 & \texttt{syn\_close\_brace} & delete a closing brace & 175 & 175 & 100\\
 & \texttt{syn\_open\_brace} & delete an opening brace & 139 & 139 & 100\\
 & \texttt{syn\_operand} & delete a comparison operand & 167 & 167 & 100\\
 & \texttt{syn\_keyword} & misspell a keyword & 176 & 176 & 100\\
\midrule
name & \texttt{name\_typo} & misspell a variable use & 208 & 208 & 100\\
 & \texttt{name\_delete\_decl} & delete a declaration & 182 & 182 & 100\\
 & \texttt{name\_redecl} & duplicate a declaration & 156 & 151 & 97\\
 & \texttt{name\_uninit} & remove an initialiser & 185 & 142 & 77\\
 & \texttt{name\_func} & misspell a function call & 173 & 128 & 74\\
 & \texttt{name\_assign\_input} & assign to an input & 189 & 189 & 100\\
\midrule
type & \texttt{type\_string\_assign} & assign a string/bool to a number & 194 & 150 & 77\\
 & \texttt{type\_bool\_assign} & assign a number to a bool & 189 & 189 & 100\\
 & \texttt{type\_arity} & change an argument count & 202 & 148 & 73\\
 & \texttt{type\_cond} & use a number as a condition & 170 & 170 & 100\\
 & \texttt{type\_return} & return the wrong type & 210 & 162 & 77\\
 & \texttt{type\_float\_int} & assign a float to an int & 181 & 181 & 100\\
\midrule
unit & \texttt{unit\_swap\_literal} & swap the unit of a literal & 225 & 212 & 94\\
 & \texttt{unit\_drop} & drop the unit of a literal & 218 & 208 & 95\\
 & \texttt{unit\_decl} & change a declared input unit & 210 & 206 & 98\\
 & \texttt{unit\_rate} & pass a wrong dose-rate unit & 195 & 86 & 44\\
 & \texttt{unit\_return} & return a wrong unit & 231 & 74 & 32\\
 & \texttt{unit\_add} & add quantities of different units & 215 & 214 & 100\\
\midrule
range & \texttt{range\_literal} & implausible threshold literal & 295 & 267 & 91\\
 & \texttt{range\_threshold} & threshold outside declared range & 321 & 273 & 85\\
 & \texttt{range\_decl} & invert a declared range & 323 & 323 & 100\\
 & \texttt{range\_const} & implausible dose constant & 307 & 137 & 45\\
\midrule
phi & \texttt{phi\_direct} & print a PHI input & 202 & 202 & 100\\
 & \texttt{phi\_alias} & export an alias of PHI & 219 & 219 & 100\\
 & \texttt{phi\_concat} & concatenate PHI into output & 222 & 222 & 100\\
 & \texttt{phi\_removeanon} & remove an anonymize() call & 219 & 116 & 53\\
 & \texttt{phi\_func} & leak PHI through a function & 241 & 241 & 100\\
\midrule
missing & \texttt{miss\_unguard} & remove a presence guard & 442 & 192 & 43\\
 & \texttt{miss\_new\_use} & use an optional value unguarded & 402 & 402 & 100\\
 & \texttt{miss\_make\_optional} & make a required input optional & 410 & 406 & 99\\
\midrule
logic & \texttt{logic\_div\_zero} & divide by constant zero & 174 & 174 & 100\\
 & \texttt{logic\_range\_zero} & divisor range includes zero & 193 & 101 & 52\\
 & \texttt{logic\_missing\_return} & delete a return & 166 & 109 & 66\\
 & \texttt{logic\_unreachable} & code after return & 160 & 112 & 70\\
 & \texttt{logic\_inverted} & invert a clinical comparison & 154 & 148 & 96\\
 & \texttt{logic\_wrong\_message} & mismatch alert text and rule & 165 & 163 & 99\\
 & \texttt{logic\_constcond} & constant if-condition & 193 & 193 & 100\\
\bottomrule
\end{longtable}

}

\section{Running the Code}

The repository accompanying this report has the following layout.

\begin{lstlisting}[language={},numbers=none,basicstyle=\ttfamily\footnotesize]
code/
  medscript/   lexer, parser, ast_nodes, semantic, ir, driver, equiv
               generator, mutations, corpus, features
  ml/          gan.py (conditional GAN), augment.py (SMOTE, jitter, ...)
  tests/       test_compiler.py (38 unit tests)
  examples/    triage_ok, phi_leak, inverted_rule, syntax_errors,
               missing_value, fold_demo (.med)
  medcheck.py  compiler + learned root-cause estimate
  run_experiments.py   make_figures.py
  results/     JSON results of every experiment
  models/      trained random-forest classifier
\end{lstlisting}

\begin{lstlisting}[language={},numbers=none,basicstyle=\ttfamily\footnotesize]
python -m unittest discover -s tests     # unit tests
python medcheck.py examples/phi_leak.med # compile + classify a program
python run_experiments.py                # all experiments (~45 min, 10 cores)
python make_figures.py                   # figures and LaTeX tables
\end{lstlisting}
